\section{Deployment}
\label{sec:deployment}

Seven containers on one host, joined by a single Docker bridge network
(\cref{fig:deployment}), each one a node. The course asks for several \emph{nodes}, and
the reference project ran on one host in the same way. What a single host does not give for
free is a network partition, a node that is alive and unreachable; \texttt{docker network
disconnect} produces it, leaving the container running, convinced it is alive, and sending
no FIN, which is the failure \texttt{docker kill} cannot show (\cref{sec:detectors}).

\paragraph{Images.} Two. \texttt{Dockerfile.erlang} builds every Erlang application with
\texttt{rebar3} and copies only the compiled \texttt{ebin} directories into the runtime
stage, pinned to \texttt{erlang:29.0.5}, the OTP on the developers' machines; the role of a
container is the \texttt{ERL\_APP} variable, so stations and coordinators are one image
with different environments. \texttt{Dockerfile.backoffice} deploys the WAR as
\texttt{ROOT} on Tomcat 10.1 together with \texttt{epmd}, because JInterface publishes the
hidden node's name on the port mapper of its own host.

\paragraph{Naming and configuration.} Every Erlang node is \texttt{vs@<hostname>}, short
names, because containers on a compose network resolve each other by service name. The three
coordinators share one configuration and differ by hostname alone: there is no
\texttt{COORD\_ID}, the bully priority being the position of a node in the sorted
\texttt{COORD\_NODES} list all of them hold, so a healthy start elects \texttt{coord3}.
Everything else is an environment variable with a development default (\cref{tab:params}),
so the project starts with no configuration at all, and \texttt{.env.demo} overrides the
leases and the site budget of \texttt{station1}.

Inside the containers every station listens on 8080 and 8081, so the application does not
know which station it is; the remapping to the published ports of \cref{fig:deployment}
(9101/9102 for drivers, 9201/9202 for charge points, 8080 for the web application) happens
only where the host has one port 8080.

\paragraph{Starting it.} From \texttt{src/deploy/}, \texttt{docker compose build} once
and \texttt{docker compose up -d} afterwards, with \texttt{.env.demo} as the env file. The
stations and the back office wait for MySQL's health check, which allows sixty retries at
five seconds because the first boot on a cold volume measured over six minutes on a laptop;
with the earlier twenty retries the other nodes refused to start against a database that
was merely slow. The schema is mounted read-only from \texttt{contracts/schema.sql}, and
the data lives in a named volume that survives \texttt{docker compose down}. The charge
points are not in the compose file: they are emulated hardware, and they run on the host
with \texttt{emulator/demo/world.sh}, which supervises one \texttt{cp.js} per connector
and restarts one that exits.

\paragraph{One network, after trying two.} A second network per site was added on 3
September on the premise that disconnecting a station from the cluster also cut its
chargers. Measuring showed the premise wrong (a published port keeps working while the
container is off the bridge, \cref{sec:station-failures}), and it was removed the same
day.

\paragraph{Decisions.} A container per node: several \texttt{-sname} nodes on one
operating system are how the election is tried in development, but nodes sharing a
filesystem, a network stack and a fate prove nothing about isolation, while \texttt{docker
kill} on a container is a real failure of a node. A three-line \texttt{start-node.sh}
instead of an OTP release: a release adds \texttt{vm.args}, \texttt{sys.config} and
profiles at the point where the risk to retire was making two nodes talk, and the script
shows exactly how a node is named, which is where things break.
